% !TEX program = xelatex
% Бүх лекцийн PDF-д хуваалцах преамбл. latex/build.sh үүнийг ашиглана.
\documentclass[12pt,a4paper,oneside]{book}
\usepackage{geometry}
\geometry{margin=1in}
\usepackage{amsmath,mathtools}
\usepackage{amsthm,amsfonts,amssymb}
\usepackage{enumitem}
\usepackage{hyperref}
\usepackage{bookmark}
\usepackage{tikz}
\usepackage{comment}
\usepackage{environ}
\usepackage{multicol}
\usepackage{pgfplots}
\pgfplotsset{compat=1.18}
\usepackage{mathrsfs}
\usepackage{caption}
\usetikzlibrary{arrows}

% Unicode & fonts (XeLaTeX)
\usepackage{fontspec}
\usepackage{unicode-math}
\defaultfontfeatures{Ligatures=TeX}
% Libertinus байвал (Mac дээр), үгүй бол CMU + Latin Modern (GitHub Actions дээр)
\IfFontExistsTF{Libertinus Serif}{%
  \setmainfont{Libertinus Serif}
  \setsansfont{Libertinus Sans}
  \setmonofont{Libertinus Mono}
  \setmathfont{Libertinus Math}
}{%
  \setmainfont{CMU Serif}
  \setsansfont{CMU Sans Serif}
  \setmonofont{CMU Typewriter Text}
  \setmathfont{Latin Modern Math}
}
\usepackage{graphicx} %зураг дуудах
\renewcommand{\chaptername}{Хичээл}
\renewcommand{\thechapter}{\arabic{chapter}}
% --- Theorem-like environments ---
\newtheorem{theorem}{Теорем}[chapter]
\newtheorem{proposition}{Өгүүлбэр}[chapter]
\newtheorem{lemma}{Лемма}[section]
\newtheorem{definition}{Тодорхойлолт}[chapter]
\newtheorem{example}{Жишээ}[chapter]
\theoremstyle{remark}
\newtheorem*{remark}{Тэмдэглэл}
\newtheorem*{prof}{Баталгаа}
\newtheorem*{note}{Санамж}
\newtheorem*{cor}{Мөрдөлгөө}

% --- Problem & Solution helpers ---
\newcounter{problem}
\newenvironment{problem}[1][]% optional title
{\refstepcounter{problem}\par\noindent\textbf{Бодлого~\theproblem.}~#1}{\par}

% --- Show/hide solutions toggle ---change the %sign
\newif\ifshowsolutions
%\showsolutionsfalse 
\showsolutionstrue 

% Robust solution environment (no nesting issues)
\NewEnviron{solution}{%
  \ifshowsolutions
    \par\smallskip\noindent\textit{Бодолт. }% 
    \BODY\par\smallskip
  \fi
}

% --- List styles ---
\setlist[itemize]{topsep=3pt, itemsep=3pt, leftmargin=1.3em}
\setlist[enumerate]{topsep=4pt, itemsep=4pt, leftmargin=1.6em}

% --- Metadata ---
\newcommand{\coursename}{Нэг хувьсагчийн функцийн онол 1В (MATH102)}
\newcommand{\lectureno}{}
\newcommand{\topic}{}
\newcommand{\semester}{2026 оны намар}
\newcommand{\smallbigcup}{\mathop{\textstyle\bigcup}}
\newcommand{\smallbigcap}{\mathop{\textstyle\bigcap}}
\newcommand{\N}{\mathbb{N}}
\newcommand{\Z}{\mathbb{Z}}
\newcommand{\Q}{\mathbb{Q}}
\newcommand{\R}{\mathbb{R}}
\newcommand{\Sx}{\mathrm{S}} % Successor
\newcommand{\itcol}[1]{{\it{\textcolor{red}{#1}}}}
\newcommand{\arccot}{\mathbf{arccot}}
\graphicspath{{../pictures/}{../}}
